% Junior Level 1, Week 1 -- answer key
\answersweek{1}{Robots Need Instructions}{Answers for workbook chapter 1 and play-at-home card 1.}

\section{Lesson questions}
\begin{answer}[What do you like to make? (warm-up chat)]
There are no right or wrong answers: the chat is ungraded. Note one or two ideas per child (for example \enquote{draws dragons}, \enquote{likes racing games}, \enquote{makes up songs}) and use them in free-build minutes and the week 9 project.
\end{answer}

\begin{answer}[Exit question]
No. The cards are the same, but the order is different, so the robot takes a different path. FORWARD, TURN RIGHT, FORWARD ends one square ahead and one square to the right. FORWARD, FORWARD, TURN RIGHT ends two squares ahead, facing right. Any answer like \enquote{the order changes where it goes} shows understanding.
\end{answer}

\begin{answer}[Reflection]
\enquote{What does a robot need from us?} Exact steps, in the right order (or \enquote{the cards}, \enquote{one step at a time}). \enquote{What do we do when an app asks for something?} Stop and ask a grown-up. Any drawing of themselves as a robot is right.
\end{answer}

\section{Workbook exercises}
\begin{answer}[Exercise 1.1: Parent Robot]
\textbf{What you should see:} a row of cards that takes the robot from START to the star, one card at a time. There are many correct rows: for a star two squares ahead and one to the right, both FORWARD, FORWARD, TURN RIGHT, FORWARD and TURN RIGHT, FORWARD, TURN LEFT, FORWARD, FORWARD work.

\textbf{What to look for:} \tickbox{} cards in one row, all pointing up \quad \tickbox{} the Helper reads one card at a time \quad \tickbox{} the children swap jobs on the timer.
\end{answer}

\begin{answer}[Exercise 1.2: The Wrong Card]
The circled card is TURN RIGHT (the second card). With the bug, the robot goes up one square, turns right too early and walks along the middle row until the third FORWARD takes it off the side of the grid. The fixed row is FORWARD, FORWARD, TURN RIGHT, FORWARD, FORWARD: up two squares to the top-left corner, turn right, then two squares to the star in the top-right corner. Some children circle the third card (FORWARD) and move it in front of TURN RIGHT: that gives the same fixed row, so accept it too.
\end{answer}

\begin{answer}[Exercise 1.3: My Own Path]
Any path is right if the robot reaches the child's star when it walks the cards. Check by walking it (with a finger on the grid or on the floor). If it misses, the child moves or changes one card and tests again: that is debugging, and it is the best possible result.
\end{answer}

\section{Play-at-home card}
\begin{answer}[Robot in the kitchen]
Look for one step at a time, said in order, and turns that are on the spot. If the child says \enquote{go to the fridge}, the robot answers \enquote{Does not compute} and the child breaks it into steps. The robot reaching the goal after a fix counts as a full success.
\end{answer}

\begin{answer}[Ask a grown-up]
The child points to any screen that asks for something (a name, a photo, an account, a purchase, an \enquote{Allow} button) and says they would stop and ask a grown-up. If the app never asks, talking about what they would do is enough.
\end{answer}
